% Section 3.3, from lectures/L03/appendix/c_button_driver.md.

\section{The Button Driver}
\label{c3:sec:driver}

\subsection{The task}
\label{c3:sec:task}
Write \file{drivers/source/btn.asm}, and extend \file{drivers/app/main.asm} with a handler.

The button driver is the LED driver's mirror image for its first half and something new for its
second. Both are handed an Arduino pin number and build a structure of register addresses; where
\code{led_init} makes the pin an output, \code{btn_init} makes it an input and switches its pull-up
on. \code{btn_pressed} is \code{led_enabled} with its answer inverted. The new part is the four
subroutines that turn a pin change interrupt on and off for one pin without disturbing the other
seven that share its vector.

\subsection{The structure}
\label{c3:sec:layout}
Ten bytes, as \file{drivers/include/btn.inc} declares.

\begin{booktable}{@{}p{3.4em}p{9em}p{2.6em}L@{}}
  \thead{Offset} & \thead{Field} & \thead{Size} & \thead{Holds} \\
  \midrule
  0 & \code{BTN_PIN_REG} & 2 & Data space address of \code{PINx} \\
  2 & \code{BTN_DIR_REG} & 2 & Data space address of \code{DDRx} \\
  4 & \code{BTN_PORT_REG} & 2 & Data space address of \code{PORTx} \\
  6 & \code{BTN_PCMSK_REG} & 2 & \code{PCMSK0} for port B, \code{PCMSK2} for port D \\
  8 & \code{BTN_PCIE_BIT} & 1 & \code{PCIE0} for port B, \code{PCIE2} for port D \\
  9 & \code{BTN_PIN} & 1 & Bit number within the port, 0 to 7 \\
    & \code{BTN_SIZE} & 10 & Total \\
\end{booktable}

The first three fields are the LED's three, at the same offsets and in the same order. That is
deliberate: a button is an LED's three register pointers plus what it takes to raise an interrupt.

It is also a trap worth naming now, because nothing prevents you passing a button structure to
\code{led_on}. \textbf{Six} bytes at the front are laid out identically, the three pointers, and
the seventh is not: offset 6 is \code{LED_PIN} in one structure and the low byte of
\code{BTN_PCMSK_REG} in the other. So \code{led_on} finds three valid port registers and then a pin
number of \code{0x6B} or \code{0x6D}, which is the same failure as reaching for the wrong field
within one structure (\secref{c3:sec:irq}): \code{shift_bits} returns 0 for any argument of 8 or
more, the mask is zero, and the subroutine writes nothing at all and reports success.

\subsection{\code{btn_init}}
\label{c3:sec:init}
\noindent\textbf{Arguments:} \code{r25:r24} is the structure address, \code{r22} the Arduino pin
number.\\
\textbf{Returns:} \code{r24} is 0 on success, 1 if the pin is not one this driver accepts.

\medskip
\noindent The same shape as \code{led_init} (\secref{c2:sec:init}), with two more fields to fill
and the direction reversed:
\begin{enumerate}
  \item Copy the structure pointer into \code{Z}.
  \item Validate, choose the port, and subtract 8 for port B. Reject before writing anything.
  \item Store the bit number, the three port register addresses, \textbf{the mask register
    address} (\code{PCMSK0} or \code{PCMSK2}) and \textbf{the \code{PCICR} bit number}
    (\code{PCIE0} or \code{PCIE2}).
  \item Make the pin an \textbf{input}: call \code{shift_bits_inverted}, then read \code{DDRx},
    \code{and} the mask, write it back.
  \item Switch the \textbf{pull-up} on: call \code{shift_bits}, then read \code{PORTx}, \code{or}
    the mask, write it back.
  \item Return 0.
\end{enumerate}

\subsubsection{Pinned details}
\textbf{Store the \code{PCICR} bit number before you call anything.} Whatever register you loaded
it into is call-clobbered, and \code{shift_bits} will take it (\secref{c2:sec:nested}). The same
applies to the mask register address. Fill the structure first, then configure the hardware.

\textbf{Steps 4 and 5 are \code{led_init}'s two steps with the masks swapped.} The LED clears
\code{PORTx} and sets \code{DDRx}; the button clears \code{DDRx} and sets \code{PORTx}. If you find
yourself writing something structurally different, one of the two is wrong.

\textbf{\code{PCMSK0} and \code{PCMSK2} are at \code{0x6B} and \code{0x6D}}, which is extended I/O,
so nothing here can use \code{in} or \code{out} (\secref{c3:sec:pcint}).

\subsection{\code{btn_pressed}}
\label{c3:sec:pressed}
\noindent\textbf{Arguments:} \code{r25:r24} is the structure address.\\
\textbf{Returns:} \code{r24} is 1 if the button is down, 0 if it is up.

\medskip
\noindent The same shape as \code{led_enabled} (\secref{c2:sec:enabled}), with one difference that
is the whole reason the subroutine exists:
\begin{enumerate}
  \item Copy the structure pointer into \code{Z}.
  \item Read the pin number from \code{BTN_PIN} and call \code{shift_bits} to turn it into a mask.
  \item Read \code{PINx} through the address in \code{BTN_PIN_REG}.
  \item Mask the two together, and return \textbf{1 when the result is zero}.
\end{enumerate}

\textbf{Step 4 is inverted, and the inversion is the point.} The pull-up holds the pin high until
the button shorts it to ground, so \code{PINx} reads 0 when the button is down
(\secref{c2:sec:button}). Every other ``is it on'' subroutine in this library returns 1 for a set
bit. This one returns 1 for a clear one, and writing it the other way round produces a program that
is exactly and reliably backwards.

\subsubsection{Pinned details}
\textbf{Read \code{PINx}, never \code{PORTx}.} On an input pin \code{PORTx} is the pull-up switch,
so reading it would report what you configured rather than what the world is doing, and it would
give the same answer pressed and unpressed. This is \code{led_enabled}'s rule with the stakes
raised: there, the two registers agree because the pin is an output; here they never agree.

\textbf{The branch costs a cycle on one path.} As in \code{led_enabled}, the natural
implementation ends in a conditional branch, so this subroutine costs one cycle more on one of its
two answers. Here it is the \textbf{unpressed} path that is dearer, because pressed is the case
that falls through. Notice it rather than avoiding it: the cross-check in
Exercise~\ref{c3:ex:crosscheck} asks you to account for that cycle.

\textbf{It reports the pin, not the press.} A mechanical button bounces, and this subroutine has no
memory and no timer, so two calls a millisecond apart during a bounce legitimately disagree
(\secref{c2:sec:button}). What it reports is the pin, now, and that is all it promises.

\subsection{The four interrupt subroutines}
\label{c3:sec:irq}
All four take the structure address in \code{r25:r24}.

\textbf{\code{btn_enable_interrupt}} sets the port's bit in \code{PCICR} and the pin's bit in the
structure's mask register. Both are read-modify-write.

\textbf{\code{btn_disable_interrupt}} clears the pin's bit in the mask register \textbf{and leaves
\code{PCICR} alone}. \code{PCICR} enables a whole port's vector and other buttons on that port may
still want it; clearing it here would silently switch off every one of them.

\textbf{\code{btn_interrupt_enabled}} returns 1 if the pin's bit in the mask register is set,
else 0.

\textbf{\code{btn_toggle_interrupt}} calls \code{btn_interrupt_enabled} and then one of the other
two.

\subsubsection{Pinned details}
\textbf{The driver does not touch the global interrupt flag.} No \code{sei}, no \code{cli},
anywhere in this file. You will find AVR drivers that run \code{sei} at the end of their enable
function, and it is convenient and it is rude: it decides on your behalf that the whole program is
now ready to be interrupted, at a moment chosen by whoever happened to configure a button last.
\code{sei} belongs in your setup code, once, after everything is configured.

\textbf{\code{btn_interrupt_enabled} must return 1 for enabled and 0 for disabled}, and that
sentence is here because the classic mistake is to return them the wrong way round. Nothing about
such code looks wrong: the branch is there, both constants are there, and they are swapped. Every
caller then does the opposite of what it meant, and \code{btn_toggle_interrupt} stops toggling and
starts latching, which is a symptom two steps removed from its cause.

\textbf{\code{btn_disable_interrupt} shifts by the pin number}, field at offset 9, not by anything
at offset 0. Reaching for the wrong field gives a shift count of \code{0x23}, which produces a mask
of zero, which clears nothing, so a button you disabled stays enabled. And clear that bit alone: a
second button on the same port has its own bit in the same mask register, and the test with two
buttons is the one that checks it survives.

\textbf{Preserve the structure pointer across the nested call in \code{toggle}.}
\code{btn_interrupt_enabled} sets \code{Z} to the same value \code{toggle} had, so \code{Z} survives
by coincidence, and relying on that is a decision rather than a fact about the contract. Pushing
\code{r25:r24} and popping them after is four instructions and eight cycles and needs no footnote.
Then jump to the enable or disable subroutine with \code{rjmp} rather than calling it: it is the
last thing you do, so its \code{ret} can be yours.

\subsection{The handler}
\label{c3:sec:handler}
Add to \file{drivers/app/main.asm}:

\textbf{A vector entry.} \code{.org} the word address of PCINT0, writing \code{.org PCI0addr} so
that the device file supplies it, then \code{rjmp} to your handler (\secref{c3:sec:table}).
AVRASM2 counts \code{.org} in words, so the datasheet's \code{0x0006} is the number you write;
\code{0x000C} is the byte address and lands you in the watchdog's slot.

\textbf{A handler} that saves what it uses, asks the button whether it is pressed, toggles the LED
if it is, restores, and executes \code{reti}. Its shape is the prologue and epilogue from
\secref{c3:sec:sreg}, with a handful of instructions in the middle. \textbf{Call it
\code{isr_pcint0}}, exactly, because two things outside your program look for it by that name: the
measurement command in the exercises, and the test suite.

\textbf{Setup} that initialises the LED on pin 13 and the button on pin 12, enables the button's
interrupt, and then runs \code{sei} \textbf{once}, after everything else. Then a main loop that
does nothing.

\subsubsection{Pinned details}
\textbf{Save \code{r24}, \code{r25}, \code{r18}, \code{r19}, \code{X} and \code{Z} as well
as \code{SREG}: nine things.} The handler passes a pointer in \code{r25:r24} and calls subroutines
that clobber both, plus \code{r18} and \code{r19} inside \code{shift_bits}, plus \code{Z}, which
\code{btn_pressed} and \code{led_toggle} load with their structure, and \code{X}, which they load
with a port register's address. Every one of those has to be saved, because the interrupted code
agreed to nothing. Leaving out the four you never name yourself is the mistake that costs sixteen
cycles less and a bug that appears only when the main loop happens to be using a pointer.

That list is built from what these drivers actually do, which is a decision
(\secref{c2:sec:nested}). The contract would let them clobber \code{r20} to \code{r23} as well,
and a handler that calls code you did not write has to save those too.

\textbf{The LED and the button must be on different pins.} Putting both on 13 means
\code{btn_init} makes the pin an input after \code{led_init} made it an output, and the LED stops
working in a way that looks like the handler never runs.

\textbf{\code{sei} goes last, in setup, once.} Not in the driver, and not before the structures
are built. With \code{sei} last nothing can interrupt setup, so the order of everything before it
stops mattering. With it first, setup is safe only while no interrupt source is switched on before
the structure its handler reads, and one that arrives early runs your handler against a structure
that is half zeros.

\textbf{This is the only part of the book a unit test cannot reach.} Everything else you have
written is a subroutine, and a test can call a subroutine; an interrupt has to arrive. So the
suite's \file{app_test.cpp} loads the whole program, runs it, and changes the pin from outside,
which is the only way to find out whether the vector entry points at your handler, whether
\code{sei} ran, and whether the handler ends with \code{reti}. It checks no cycle counts: what your
handler costs depends on what your handler saves, and measuring that is
Exercise~\ref{c3:ex:crosscheck}'s job rather than the suite's.

\subsection{What this driver does not do}
\label{c3:sec:notdo}
\textbf{It does not debounce.} A press produces several interrupts, a few milliseconds apart, and
the handler runs for every one. An LED toggled by this driver will appear to respond to about half
of your presses, and that is the driver working correctly.

\textbf{It does not tell you which way the pin went.} The interrupt fires on press and on release,
and by the time the handler reads the pin it is reading the present, not the edge. Comparing
against a remembered previous state is the usual answer, and it needs a variable shared with the
handler (\secref{c3:sec:shared}).

\textbf{It does not tell you which pin.} One vector serves eight, so a handler with two buttons on
one port has to ask both.

\textbf{It does not stop you enabling an interrupt with no handler behind it.} The vector then
contains whatever was there, which for an unwritten slot is \code{0xFFFF}, and the machine executes
off into unprogrammed flash and eventually restarts (Exercise~\ref{c1:ex:program}, part d).
